\section{Introduction}

Strangeness enhancement is considered as one of the first proposed key signatures \cite{Muller:1982} for the formation of the quark--gluon plasma 
(QGP), a strongly interacting quantum chromodynamics state of nuclear matter with deconfined quarks and gluons. In ultrarelativistic heavy-ion 
collisions with sufficiently high temperatures and energy densities where QGP formation is anticipated, an enhanced production of strange hadrons 
with respect to small systems (like pp or p--Be collisions) is observed \cite{WA97:1999uwz, STAR:2008, PbPb276:2014}, with an emphasis on relative 
yield of multi-strange baryons to non-strange hadrons being more pronounced.
In smaller collision systems such as proton--proton (pp) and proton--lead (p--Pb), strangeness enhancement can be described as an increasing ratio 
of strange hadron yields over pion yields as a function of charged particle multiplicity. Multiplicity dependent studies performed at the ALICE 
experiment at the LHC have demonstrated a smooth evolution of this ratio across multiple collision systems with different collision energies. 
The increase is more prominent in multi-strange particles, indicating that the enhancement scales with particle strangeness content \cite{Nature:2017}. 